%% kapa: exam=1 round=27 year=2016 period=1 form=A subject=civil_law subject_ko=민법 seq=26 number=26
%% source: Q-Net 감정평가사 기출문제 (cst00309) artlSeq=5169402 file=extracted/1008/5169402/제27회 감정평가사 1차 문제/제27회 감정평가사 1교시문제(A형).pdf
%% conversion: kordoc 4.24.51 -> tools/gen1.py (auto); review=proofread-sonnet; flags=none
\begin{question}{26}
2015년 5월경 명의신탁자 乙과 명의수탁자 丙의 약정에 따라, 丙은 매수인으로서 부동산 매도인 甲과 매매계약을 체결하고 대금을 지급한 후, 자신의 명의로 소유권이전등기를 경료받았다. 이에 관한 설명으로 옳은 것을 모두 고른 것은? (단, 「부동산 실권리자명의 등기에 관한 법률」 제8조(종중, 배우자 및 종교단체에 관한 특례) 등에 해당하는 예외 사유가 없으며, 다툼이 있으면 판례에 따름)
\begin{framedtext}
ㄱ. 甲이 명의신탁약정을 알지 못한 경우, 乙은 丙에 대하여 소유권이전을 청구할 수 있다. \par
ㄴ. 甲이 명의신탁약정을 알고 있었던 경우, 丙은 甲에 대하여 매매대금의 반환을 청구할 수 있다. \par
ㄷ. 甲이 명의신탁약정을 알고 있었던 경우, 乙은 甲에 대하여 진정명의회복을 원인으로 한 소유권이전등기를 청구할 수 있다. \par
ㄹ. 甲이 명의신탁약정을 알고 있었던 경우, 乙은 甲을 대위하여 丙 명의 등기의 말소를 청구함과 동시에 甲에 대하여 매수인의 지위에서 소유권이전등기를 청구할 수 있다. \par
\end{framedtext}
\begin{choices}
  \item ㄱ
  \item ㄴ
  \item ㄴ, ㄹ
  \item ㄷ, ㄹ
  \item ㄱ, ㄷ, ㄹ
\end{choices}
\end{question}
\answer{2}
